\SourceLeaf{10}
\begin{center}
{\Large\textsc{PREFACE:}}

\medskip
\BookRule{0.07\linewidth}
\end{center}

The investigations in this work belong to the part of mathematics that deals with integers; fractions are mostly excluded, and irrational quantities always are. What is called indeterminate or Diophantine analysis, which teaches us to select, from the infinitely many solutions of an indeterminate problem, those expressed by integers or at least by rational numbers (usually with the further condition that they be positive), is not this discipline itself but a very special part of it. Its relation to the discipline is much like that of the art of reducing and solving equations (Algebra) to the whole of Analysis. Just as all investigations into the general properties of quantities belong to the domain of \emph{Analysis}, so integers (and fractions insofar as they are determined through integers) are the proper subject of \textsc{Arithmetic}. But what is usually taught under the name of Arithmetic hardly goes beyond the art of writing numbers and calculating (that is, representing numbers by suitable signs, for example in the decimal system, and carrying out arithmetical operations). Some further matters are added that either do not belong to Arithmetic at all (such as the theory of logarithms), or at least are not specific to integers but apply to all quantities. So it seems useful to distinguish two parts of Arithmetic: to assign these matters to elementary Arithmetic, and to reserve all general investigations into the properties specific to integers for \emph{Higher Arithmetic}, which alone concerns us here.

Higher Arithmetic includes the matters that Euclid presented in Book VII and the following books of the Elements, with the elegance and rigour usual among the ancients; they are, however, confined to the very beginnings of this science. The celebrated work of Diophantus, devoted entirely to \SourceLeaf{11}indeterminate problems, contains many questions whose difficulty and subtle methods inspire considerable respect for the author's ingenuity and insight, especially when we consider how few resources were available to him. But these problems call for dexterity and skilful handling rather than deeper principles; they are also too special and rarely lead to more general conclusions. So this book seems to mark an epoch in the history of mathematics more because it contains the first traces of symbolic notation and Algebra than because it enriched Higher Arithmetic with new discoveries. By far the greater part is owed to more recent scholars. Among them a few men, but men of immortal glory---P. \textsc{de Fermat}, L. \textsc{Euler}, L. \textsc{La Grange}, A. M. \textsc{Le Gendre} (to pass over a few others)---opened the way into the inner sanctuary of this divine science and revealed the riches it holds. I will not recount here which discoveries came from each of these mathematicians. These can be learned from the prefaces to the Additions with which the distinguished La Grange enriched Euler's Algebra, and to the work recently published by the distinguished Le Gendre that I will mention shortly. Most are also cited in their proper places in these Arithmetical Investigations.

The purpose of this work, which I publicly promised to publish five years ago, was to make known my investigations in Higher Arithmetic, some undertaken before that time and some afterwards. To explain why the science is taken up here almost from its beginnings, and why many investigations to which others have already devoted their efforts are taken up again, I thought I should point out that, when I first turned to this kind of investigation at the beginning of 1795, I knew nothing of what more recent scholars had accomplished in this field and had no means of learning about it. While occupied with another task, I happened upon a remarkable arithmetical truth (unless I am mistaken, it was the theorem of article 108). I found it very beautiful in itself and suspected that it was connected with greater things, so I worked as hard as I could to understand its principles and find a rigorous proof. When I finally succeeded as I had hoped, these questions had so caught my interest that I could not leave them. As one thing kept opening the way to another, most of what is presented in the first four Sections of this work was completed before I had seen anything of the similar work of other \SourceLeaf{12}mathematicians. Later, when I had the opportunity to read the writings of these great minds, I recognised that most of my thinking had been spent on matters long since settled. But following in their footsteps, I worked all the more eagerly to develop Arithmetic further; this led to various investigations, some of which are presented in Sections V, VI and VII. When, after some time, I decided to publish the fruits of my nightly work, I let myself be persuaded---all the more readily because several people wished it---to keep even the earlier investigations. At that time there was no book from which one could learn the work of other mathematicians on these subjects, scattered through the Transactions of Academies; many of my investigations were entirely new and most used new methods; and, finally, all were so \emph{closely} connected with one another and with the later investigations that even the new material could not be explained conveniently enough without taking up the rest from the beginning.

Meanwhile an excellent work appeared by a man who had already done much for Higher Arithmetic: \emph{Le Gendre, Essai d'une théorie des nombres, Paris, year VI}. He not only carefully collected and arranged everything accomplished so far in this science, but also added a great many new discoveries of his own. This book reached me late, after most of my work had already been printed, so I could not mention it wherever the similarity of the subjects might have provided an opportunity. It seemed necessary to add in the Additions some observations on just a few passages; I hope that this most courteous and fair-minded man will take them kindly.

During the printing of this work, repeatedly interrupted and prolonged by various obstacles into a fourth year, I continued investigations I had already begun but had decided to publish later to keep the book from becoming too large, and also began several new ones. Several others that, for the same reason, I had only touched on briefly, since a fuller treatment seemed less necessary (for example, those in articles 37, 82 and following, and elsewhere), were later taken up again and led to more general investigations that seem well worth publishing (compare also what is said in the Additions about article 306). Finally, the book grew far larger than I had expected, especially because of the length of Section V, so several things originally intended for it had to be \SourceLeaf{13}cut, including the whole \emph{eighth} Section (already mentioned at various places in this volume, and containing a general treatment of algebraic congruences of any degree). All this will easily fill a volume as large as the present one, and will be published as soon as an opportunity arises.

In several difficult questions I have used synthetic proofs and left out the analysis by which they were found. The main reason was the need for brevity, which I had to consider as far as possible.

The theory of the division of the circle, or of regular polygons, treated in Section VII, does not \emph{itself}, \emph{in its own right}, belong to Arithmetic; its \emph{principles}, however, must be drawn solely from Higher Arithmetic. This may be as unexpected to mathematicians as I hope the new truths drawn from this source will be welcome to them.

These are the matters I wished to explain to the reader in advance. It is not for me to judge the work itself. My greatest wish is that what fills gaps that have remained until now, or opens the way to new discoveries, may please those who have the advancement of the sciences at heart.

\BookRule{0.16\linewidth}
